\chapter{Introduction}

\section{Background}

In the current digital era, the development of information technology has significantly transformed the employment landscape, particularly through the emergence of the gig economy. The gig economy represents a task-based, temporary work model facilitated by digital platforms, such as food delivery, online transportation, or creative freelancing. In Indonesia, this sector is growing rapidly, with the number of non-traditional workers reaching 41.6 million people in 2025, encompassing gig workers and part-timers, especially among young people including students who have side gigs or part-time jobs \cite{bps2025tenagakerja, ilo2025worldemployment}.

However, behind its flexibility, this work model often triggers mental health issues, particularly burnout---commonly known as chronic fatigue---due to income uncertainty, irregular schedules, and algorithmic pressure from platforms \cite{frontiers2025gigwork}. Studies show that around 50--60\% of gig workers experience burnout symptoms, such as emotional exhaustion and decreased performance, with a higher risk compared to permanent employees \cite{jmir2025digitalinterventions}. In Indonesia, this condition is further exacerbated by limited access to conventional mental health services, where only 0.3 psychiatrists are available per 100,000 population \cite{kemenkes2025akseskesehatan}.

Students who work part-time or have side gigs are also included in this non-traditional worker group and are likely to face a double burden between studying and working, leading to high burnout rates reaching 91\% in some local samples \cite{syiahkuala2024burnout}. In addition, the use of mental health applications based on chatbots, such as mood tracking for self-help, is increasingly popular as a more accessible alternative. However, most existing applications still use generic interfaces that are less adaptive to the specific needs of users with busy schedules and high cognitive load, such as quick input or personalization based on daily stress patterns \cite{pmc2025scopingreview}.

This gap becomes increasingly evident when mental health applications fail to meet the needs of non-traditional workers, where intuitive and personalized digital interventions are required to support daily burnout management. Therefore, there is a need for adaptive UI/UX design in a mood tracker application that can adjust to the needs of these users, thus encouraging this research to fill that gap.

\section{Research Problem Statement}

Based on the background above, the research problem statements of this study are as follows:

\begin{enumerate}
	\item The lack of adaptive UI/UX in chatbot-based mood tracker applications that support burnout management for non-traditional workers, such as students with side gigs or part-time jobs.
	\item The suboptimal interface of mental health applications in accommodating the needs of non-traditional users, such as quick input features and personalization to reduce cognitive load and enhance emotional engagement.
	\item The absence of evaluations on the effectiveness of mood tracker applications' UI/UX using relevant usability metrics tailored to the context of non-traditional workers.
\end{enumerate}

\section{Research Limitations}

To maintain focus and account for resource constraints, the limitations of this study are as follows:

\begin{enumerate}
	\item The research focuses on non-traditional workers among students who are currently pursuing academic studies while also having experiences with side gigs or part-time jobs as a representative subset of the gig economy, thus excluding full-time gig workers such as professional online motorcycle taxi drivers or freelancers on certain platforms who work more than 56 hours per week.
	\item The scope of the intervention is limited to adaptive UI/UX for mood trackers with basic features (emotion logging, simple trend visualization, and integration with prebuilt chatbots), without developing complex AI backend or clinical diagnosis capabilities.
	\item The evaluation covers aspects of usability and emotional impact, rather than long-term therapeutic effectiveness or medical validation.
\end{enumerate}

\section{Research Benefits}

This study is expected to provide the following benefits:

\begin{itemize}
	\item \textbf{Theoretical Benefits:} Contribute to the field of Human-Computer Interaction (HCI) by proposing an adaptive UI/UX model for mental health applications targeting non-traditional user groups, along with references for relevant usability evaluation metrics.
	\item \textbf{Practical Benefits:} Produce a prototype mood tracker application that can be used by students or part-time workers for daily burnout management, featuring a more intuitive and personalized interface.
	\item \textbf{Social Benefits:} Support increased mental health awareness among young non-traditional workers in Indonesia, aligning with national policies on the welfare of informal workers.
\end{itemize}

\section{Thesis Structure}

This research is organized into five chapters with the following structure:

\begin{description}
	\item[\textbf{Chapter 1: Introduction}] Contains the background of the problem, research problem statement, research limitations, research benefits, and thesis structure.
	\item[\textbf{Chapter 2: Literature Review}] Discusses foundational theories, concepts related to the gig economy, burnout, adaptive UI/UX, mood trackers, and relevant prior studies.
	\item[\textbf{Chapter 3: Research Methodology}] Explains the research design, research stages, population and sample, data collection methods, research instruments, research ethics, and data analysis methods.
	\item[\textbf{Chapter 4: Results and Discussion}] Presents the prototype testing results, analysis of data metrics (TSR-CC, ERI, PEI), and discussion of findings.
	\item[\textbf{Chapter 5: Conclusion}] Contains the research conclusions, suggestions for future development, and practical implications.
\end{description}
